\documentclass[12pt]{article}


% ============================
% Codificación y lenguaje
% ============================
\usepackage[utf8]{inputenc}
\usepackage[T1]{fontenc}
\usepackage[spanish]{babel}

% ============================
% Colores (cargar temprano para evitar choques)
% ============================
\usepackage[dvipsnames,table,xcdraw]{xcolor}

% ============================
% Márgenes y alineación
% ============================
\usepackage[left=.75in, right=.8in, top=1in, bottom=1in]{geometry}
\usepackage{ragged2e} % \justifying, etc.
\setlength{\headheight}{14.49998pt}
\setlength{\parindent}{0pt}
\setlength{\parskip}{0.75em}
% ============================
% Matemáticas
% ============================
\usepackage{amsmath, amssymb, amsfonts, amsthm}
\usepackage{mathrsfs}
\usepackage{soul}

\newtheoremstyle{caso_style}
  {3pt}           % Espacio vertical arriba
  {3pt}           % Espacio vertical abajo
  {\normalfont}   % Fuente del cuerpo del texto (usa \itshape si prefieres cursiva)
  {}              % Sangría (vacío = sin sangría)
  {\bfseries}     % Fuente del encabezado (Negrita)
  {.}             % Puntuación al final del encabezado (un punto)
  {0.5em}         % Espacio horizontal después del encabezado
  {\thmname{#1}\thmnote{ #3}}

\theoremstyle{caso_style}
\newtheorem*{caso}{Caso}


% ============================
% Gráficos y diagramas
% ============================
\usepackage{graphicx}
\usepackage{tikz-cd}
\usepackage{tikz}
\usetikzlibrary{automata, arrows.meta, positioning}

% ============================
% Estilo y utilidades
% ============================
\usepackage{fancyhdr}
\usepackage{booktabs}
\usepackage{float}
\usepackage{enumitem}
\usepackage{multicol}
\usepackage{tcolorbox}
\usepackage{pifont}
\usepackage{bbding}
\usepackage{comment}
% \usepackage{url} % NO es necesario si cargas hyperref (lo hace internamente)

% ============================
% Bibliografía APA con biblatex
% ============================
%\usepackage{csquotes} 
%\usepackage[backend=biber,style=ieee]{biblatex} % puedes cambiar ieee por apa, numeric, etc.
%\addbibresource{referencias.bib}

% ============================
% Algoritmos (una sola vez)
% ============================
\usepackage{algorithm}
\usepackage{algpseudocode}
\renewcommand{\thealgorithm}{} % Sin numeración (opcional)

% ============================
% Definición de colores personalizados
% ============================
\definecolor{mygreen}{rgb}{0.855, 0.969, 0.651}
\definecolor{myblue}{rgb}{0.68, 0.85, 0.9}
\definecolor{mythistle}{rgb}{0.847,0.749,0.847}
\definecolor{myorange}{rgb}{1.0, 0.6, 0.2}
\definecolor{mypurple}{rgb}{0.7, 0.2, 0.9}
\definecolor{myyellow}{rgb}{1.0, 1.0, 0.6}
\definecolor{myred}{rgb}{0.9, 0.3, 0.3}
\definecolor{mycyan}{rgb}{0.6, 1.0, 1.0}
\definecolor{mybrown}{rgb}{0.6, 0.3, 0.2}
\definecolor{mygray}{rgb}{0.85, 0.85, 0.85}
\definecolor{mymagenta}{rgb}{1.0, 0.5, 1.0}

% ============================
% tcolorbox: caja para pseudocódigo (tu estilo)
% ============================
\newtcolorbox{pseudocode}{
  colback=white, colframe=black, fontupper=\ttfamily,
  boxrule=0.5mm, arc=3mm,
  left=1mm, right=1mm, top=1mm, bottom=1mm
}

% ============================
% Listado de código (listings) con colores tipo terminal
% ============================
\definecolor{terminalbackground}{RGB}{0,0,0}
\definecolor{terminaltext}{RGB}{255,255,255}
\definecolor{terminalrule}{RGB}{0,128,0}
\definecolor{terminalnumber}{RGB}{128,128,128}

\usepackage{listings}
\lstdefinestyle{terminal}{
  backgroundcolor=\color{terminalbackground},
  basicstyle=\color{terminaltext}\large\bfseries\ttfamily,
  frame=trbl,
  framesep=2pt,
  framerule=0pt,
  xleftmargin=10pt,
  xrightmargin=10pt,
  rulecolor=\color{terminalrule},
  numberstyle=\tiny\color{terminalnumber},
  showstringspaces=false,
  upquote=true
}

% ============================
% Hipervínculos (cargar casi al final)
% ============================
\usepackage{hyperref}
\hypersetup{
  colorlinks=true,
  linkcolor=Orchid,
  filecolor=magenta,
  urlcolor=Orchid,
  citecolor=Orchid,
}

% ============================
% Figuras/tablas: títulos
% ============================
\usepackage{caption}
\setlength{\headheight}{14.49998pt}
\setlength{\parindent}{0pt}
\setlength{\parskip}{0.5em}
\begin{document}

\input{Portada.tex}
\clearpage
\pagestyle{fancy}
\renewcommand{\headrulewidth}{0.4pt} % línea superior
\renewcommand{\footrulewidth}{0.4pt} % línea inferior

\lhead{\scshape{\footnotesize Reporte}}
\rhead{\footnotesize{\scshape Práctica 03}}
\lfoot{\scshape{\footnotesize Facultad de Ciencias, 2027-I}}
\rfoot{\scshape{\footnotesize Compiladores}}
\cfoot{\thepage}

\section{Objetivo}

El objetico de esta práctica es el conocer e implementar un analizador sintáctico descendente
recursivo, es decir, un parser para el lenguaje MiniC que tiene como rol el consumir
la secuencia de tokens generados por el analizador léxico, verificando si los elementos 
se encuentran correctos con respecto a las producciones de la gramática.

Su papel principal es el conectar el analizador léxico con la construcción futura del AST,
que ayuda al validar el orden y la agrupación de los tokens cumpla con la gramatica
para posteriormente poder construir su representación.

\section{Evolución respecto de la Práctica 2}

\subsection{Módulos modificados o agregados}

Para esto, se incorpora los componentes del analizador sintáctico por medio de la carpeta de
\texttt{parser/} dentro de la estructura del proyecto, generando los archivos \texttt{parser.h, 
parser.c}. De este modo, el archivo \texttt{main.c} fue ampliado para que se pueda
coordinar el lexer y el parser del compilador.

\subsection{Interfaz incremental incorporada al lexer}

Para este momento del proyecto, se cambia la funcionalidad del analizador léxico para dejar de procesar
la entrada por bloque y se usa una interfaz incremental, donde cada vez que el parser realiza
una lectura del token del archivo, el lexer procesa ese único token disponible para conservar el estado
interno y de ese punto continuar con la siguiente llamada del parser.

\subsection{Eliminación de la impresión de tokens como efecto secundario}

Durante la práctica pasada, el lexer imprmía los tokens en la consola para su verificación
pero en estos momentos, la salida fue eliminada para que los tokens sean consumidos
como entrada del parser.

\subsection{Refactorizaciones realizadas}

Para este caso, las principales refactorizaciones del proyecto fue en la implementación de 
\texttt{lexer.c}, al agregar una nueva interfaz en el \texttt{lexer.h} para indicar
el estado del lexer, si hay error en la entrada o salida del programa y si hay errores
de memoria. Además de agregar una estructura para el lexer, dejando el archivo 
a leer, la columna y fila actuales, y si se encontro el final del archivo.

A su vez, el ciclo principal del lexer se paso en su nueva función de \texttt{lexer\_next\_token}
como la destrucción del lexer para la limpieza en memoria.

\subsection{Funcionalidades léxicas conservadas}

El sistema principal de nuestro lexer se mantuvo intacto ante las nuevas exigencias de esta 
práctica, el manejo y lectura del archivo fuente, el sistema de almacenamiento de tokens,
el seguimiento de línea y columna, la gestión de espacios en blanco y comentarios, 
el reporte de errores léxicos, el reconocimiento de identificadores, palabras resevadas,
literales y operadores.

\section{Arquitectura e interfaz}

Para esto, el analizador sintáctico está estrcuturado bajo el paradigma de descenso recursivo,
teniendo su lógica en funciones individuales que son traducciones de las producciones,
la cual se encuentra encapsulada en su propio módulos y con una cabecera pública que expone
sus funciones principales y la estructura de estado del parser.

En esta misma cabecera, se dejo lista la estructura para representar el estado actuales
del parser que mantiene un puntero hacia el analizador léxico instanciado, teniendo banderas de control
como la gestión de errores (como si hay un error o el modo pánico), el token actual y el anterior
teniendo como función la anticipación de tokens para las decisiones futuras del parser.

De esta manera, podemos definir las responsabilidades del lexer, parser y \texttt{main.c}
como:

\begin{itemize}
    \item El lexer tiene el rol de leer el archivo y proporcionar tokens válidos
    o inválidos bajo las peticiones del parser, sin el almacenamiento de los mismos.
    \item El parser se encarga de solicitar los tokens, verificar que su secuencia cumpla
    con la gramática, reportar fallas y tener un mecanismo de sincronización ante errores.
    \item Por último, el main es el orquestador de ambas fases, validando los argumentos
    de entrada, inicializando el lexer y el parser, hace el análisis global, reporta el resultado
    y libera los recursos antes de la salida del programa.
\end{itemize}

Con respecto a la liberación de memeoria, durante la ejecución del parser al solicitar
tokens al lexer, adquiere el manejo absoluto de la memoria asociada al lexema, lo cual durante su 
avance el parser se encarga de liberar la memoria del token anterior llamando la 
función de 

\verb|token_destroy(&parser->previous)| antes de escribir el actual. Por último
del programa se llama la función \verb|parser_destroy| auxiliar para limpiar cualquier token remanente del estado.

Finalmente, para tratar los fallos internos del lexer, se utiliza la estructura de estado del 
lexer, ya que nuestra función de \verb|parser_advance| intercepta este estado y detiene el ciclo de advance
para informar en la salida de errores. Además de que, el analizador léxico puede devolver
el token de error y en este caso, el parser imprime el diagnóstico léxico,
marca como inválido, destruye el token y continúa al siguiente token válido.

\section{Descenso recursivo y gramática}

La técnica que se usa, establece que el mapeo directo de un uno a uno entre
los símbolos no terminales de nuestra gramática, la cual fue mapeada directamente a su función
dentro del parser, para validar su propia estrcutura mediante el consumo de los tokens
y la invoación de otras producciones.

De este modo, nuestra funcion para revisar la regla gramatical dentro de \verb|parse_statement|
tiene que inspeccionar el token actual, de ahí se usa un switch para enrutar el control hacia alguna función
adecuada siguiendo las reglas gramaticales.

\section{Expresiones} 

La gramática de expresiones se implementa mediante descenso recursivo por niveles, 
donde cada nivel de precedencia cuenta con su propia función. El análisis inicia 
en el nivel de menor prioridad y escala por medio de llamadas secuenciales hacia los 
niveles superiores, verificando que las operaciones con mayor precedencia se agrupen 
primero. La asociatividad por la izquierda para los operadores binarios se maneja 
iterando con bucles while, mientras que la asociatividad por la derecha del 
operador unario negativo se logra mediante recursión directa sobre su propia rama. 
Finalmente, la función base \verb|parse_primary| consume los simbolos terminales 
y procesa las agrupaciones, generando un nuevo ciclo de evaluaciones,

\section{Diagnósticos y recuperación}

Para el diagnostico de los errores sintacticos es por medio de emitir cuando el 
token de anticipación no coincide con la estructura gramatical esperada
 por la regla actual, imprimiendo la posición , el elemento esperado y el lexema 
 encontrado. 
 
Se diferencia de un token léxico ERROR en que este último proviene de un 
fallo de escaneo del lexer y es consumido de inmediato registrando la falla léxica, 
evitando emitir un error sintáctico redundante. 

Para recuperarse de un fallo, el parser activa el modo pánico y sincroniza avanzando 
hasta encontrar puntos seguros como delimitadores, inicios de sentencia o el final del 
archivo . El progreso del análisis y la prevención de ciclos 
infinitos se garantizan forzando el avance sobre tokens no reconocidos o 
devolviendo el control a la estructura superior en límites de bloque.

Finalmente, la validación completa del archivo se verifica en \verb|parse_program|,
 la cual exige el consumo explícito de \verb|TOKEN_EOF| tras procesar la lista de 
 sentencias para asegurar que no quede contenido extra sin analizar.

\section{Resultados y pruebas}

Para las pruebas, se ejecutaron los tests públicos del proyecto para esta fase del compilador
donde se pudo observar que todos pasaron sin problemas. Además de que se agregaron als pruebas
de regresión de las prácticas pasadas, dejandolas aparte de su ejecución en otra regla 
del archivo \texttt{Makefile}.

A continuación se puede observar los logs de ambas pruebas:

\begin{lstlisting}[language=BASH, breaklines=true]
> make test                                                      
[OK] p01_empty
[OK] p02_declarations
[OK] p03_assignment_print
[OK] p04_expression_precedence
[OK] p05_if_else
[OK] p06_dangling_else
[OK] p07_while_blocks
[OK] p08_complete_program
[OK] p09_missing_semicolon
[OK] p10_missing_parenthesis
[OK] p11_missing_brace
[OK] p12_incomplete_expression
[OK] p13_operator_without_operand
[OK] p14_unexpected_token
[OK] p17_trailing_content
[OK] p15_multiple_errors
[OK] p16_lexical_error

Aprobadas: 17
Fallidas: 0

> make test-regression
[OK] p01_identifiers
[OK] p02_reserved_words
[OK] p03_case_sensitive
[OK] p04_booleans
[OK] p05_compound_operators
[OK] p06_incomplete_operators
[OK] p07_comments
[OK] p08_comment_at_eof
[OK] p09_slash_and_comments
[OK] p10_longest_match
[OK] p10_mixed_basic
[OK] p11_error_recovery
[OK] p12_complete_lexer
[OK] p13_crlf
[OK] p1_empty
[OK] p2_whitespace
[OK] p3_single_symbols
[OK] p4_integer
[OK] p5_integer_expression
[OK] p6_multiple_lines
[OK] p7_tabs_spaces
[OK] p8_error_character
[OK] p9_negative_number

Aprobadas: 23
Fallidas: 0

\end{lstlisting}

Además de que se revisó el uso de la memoria con el programa \texttt{valgrind}, checando
cuanta memoria se ocupó en la ejecución de nuestro compilador y cuanta memoria se liberó,
indicando si hubo fuga de memoria. En este caso, se realizó para la prueba con
nombre \verb|./tests/public/inputs/p15_multiple_errors.mc| y se obtuvo lo siguiente:

\begin{lstlisting}[language=BASH, breaklines=true]
> make memcheck FILE=./tests/public/inputs/p15_multiple_errors.mc
==144845== Memcheck, a memory error detector
==144845== Copyright (C) 2002-2024, and GNU GPL'd, by Julian Seward et al.
==144845== Using Valgrind-3.25.1 and LibVEX; rerun with -h for copyright info
==144845== Command: ./minic ./tests/public/inputs/p15_multiple_errors.mc
==144845== 
Error sintactico [2:0]: se esperaba ';' al final de la declaracion, pero se encontro 'print'.
Error sintactico [4:4]: se esperaba una expresion, pero se encontro ';'.
==144845== 
==144845== HEAP SUMMARY:
==144845==     in use at exit: 0 bytes in 0 blocks
==144845==   total heap usage: 26 allocs, 26 frees, 4,708 bytes allocated
==144845== 
==144845== All heap blocks were freed -- no leaks are possible
==144845== 
==144845== For lists of detected and suppressed errors, rerun with: -s
==144845== ERROR SUMMARY: 0 errors from 0 contexts (suppressed: 0 from 0)
\end{lstlisting}

\section{Uso de herramientas de Inteligencia Artificial}

No se utilizo ninguna herramienta de inteligencia artificial para la elaboración
de la práctica.

\section{Conclusiones}

El desarrollo del analizador sintáctico descendente recursivo de MiniC 
requirió superar desafíos como ciclos infinitos en el procesamiento, 
el silenciamiento incorrecto de diagnósticos múltiples y la necesidad de 
adaptar el entorno de pruebas para capturar la salida de errores. 
La transición hacia una interfaz incremental en el lexer optimizó 
el uso de memoria, mientras que la técnica de descenso recursivo 
por niveles demostró ser altamente efectiva para validar la gramática 
y gestionar correctamente la precedencia y asociatividad de las expresiones. 
Finalmente, la exitosa implementación del modo pánico ayudo al compilador en seguir
ejecutando para recuperarse de errores y continuar el análisis, 
consolidando una arquitectura para una futura implementación de AST  

\section{Referencias}
\begin{itemize}
    \item Pablo Martínez, Y. J.; Salazar González, P. Y. (2026). Compiladores - Práctica 3, semestre 2027-1. Facultad de Ciencias, UNAM.
\end{itemize}
\end{document}
